\documentclass{article}

\usepackage{amsmath,amssymb}
\usepackage{xurl}
\usepackage[hidelinks]{hyperref}
\setlength{\emergencystretch}{2em}

\input{command}

\title{Exact-Rank Wording in Wolf's Choi and Witness Criteria}
\author{}
\date{2026-08-24}

\begin{document}
\maketitle
\gapnote{false-source}{resolved}

\section{The printed criterion}

Let $T:M_d(\C)\to M_{d'}(\C)$ be linear and let
\begin{align}
    \tau_T
      &=(T\otimes\operatorname{id}_d)
        \bigl(\ketbra{\Omega_d}{\Omega_d}\bigr)
      \in M_{d'}(\C)\otimes M_d(\C).
\end{align}
Wolf's Proposition~3.1 assumes $k\le d$ and states, among its three
equivalent conditions, that
\begin{align}
    \langle\psi,\tau_T\psi\rangle&\ge0
    &&\text{whenever }\operatorname{SR}(\psi)=k.
    \tag{Wolf 3.1(3)}
\end{align}
The local source prints ``Schmidt-rank $n$'' both in the statement and in the
last sentence of the proof.
\footnote{Source: \path{Notes/WolfNoteTexSource/ch03_positive_not_completely.tex},
lines 89--115. Tracked by
\href{https://github.com/LionSR/QICLean/issues/24}{QICLean issue \#24}.}

\section{The dimension defect}

Every vector in $\C^{d'}\otimes\C^d$ satisfies
\begin{align}
    \operatorname{SR}(\psi)&\le\min\{d',d\}.
\end{align}
Consequently, when $d'<k\le d$, there is no vector of Schmidt rank exactly
$k$. Condition~(Wolf 3.1(3)) is then vacuous and cannot characterize
$k$-positivity.

For example, take $d=2$, $d'=1$, $k=2$, and
$T(X)=-\tr(X)$. There is no Schmidt-rank-two vector in
$\C^1\otimes\C^2$, so the printed quadratic-form condition holds. On the
other hand, applying $T\otimes\operatorname{id}_2$ to a nonzero product
projector gives a negative rank-one matrix. Thus $T$ is not $2$-positive.

The last step of the printed proof has a second, related loss: a vector of the
form $(\one\otimes X)^\dagger\widetilde\psi$ has Schmidt rank at most
$\operatorname{rank}(X)$, but need not have rank equal to
$\operatorname{rank}(X)$. Hence that step
directly proves the bounded-rank condition, not the exact-rank condition.

\section{The corrected statement}

The dimension-generic Choi criterion is
\begin{align}
    T\text{ is $k$-positive}
    \quad\Longleftrightarrow\quad
    \langle\psi,\tau_T\psi\rangle&\ge0
    &&\text{for every }\psi\text{ with }\operatorname{SR}(\psi)\le k.
    \tag{corrected Wolf 3.1(3)}
\end{align}
This is the condition proved by the compression argument: factoring the
coefficient matrix of $\psi$ through $\C^k$ writes
$\psi=R_X^\dagger\eta$, and positivity of the right compression
$R_X\tau_T R_X^\dagger$ gives the desired quadratic-form inequality.

For $d>0$, the formal theorem covers independent input and output dimensions,
including $k=0$ and $d'=0$.
\footnote{Formal declarations:
\leanid{ChoiJamiolkowski.isNPositiveMap_iff_forall_hasSchmidtRankLE_choiMatrix_quadraticForm_nonneg_rectangular}
and
\leanid{ChoiJamiolkowski.isNPositiveMap_iff_forall_rankProjection_choiMatrix_sandwich_posSemidef_rectangular}
in \doclink{QICLean/Channel/Schwarz/ChoiCompression}.}
The projection formulation additionally assumes $k\le d$, exactly as needed
to choose a rank-$k$ projection on the input factor.

If one also assumes $k\le d'$, an exact-rank formulation may be recovered
from the bounded-rank formulation by a separate density argument: coefficient
matrices of rank exactly $k$ are dense among matrices of rank at most $k$, and
the quadratic form is continuous. This argument is not the proof printed by
Wolf and is not used in the formal criterion.

\section{The same defect in Proposition~3.3}

Wolf's Proposition~3.3 considers a density operator
$\rho\in M_d(\C)\otimes M_{d'}(\C)$ and states that its Schmidt number is at
least $n+1$ if and only if there is a Hermitian $W$ with
$\tr(W\rho)<0$ and
\begin{align}
    \langle\psi,W\psi\rangle&\ge0
    &&\text{for every $\psi$ with $\operatorname{SR}(\psi)=n$}.
    \tag{Wolf 3.3}
\end{align}
The source then defines $S_n$ as the density operators of Schmidt number at
most $n$ and separates $\rho$ from $S_n$; its converse says that $S_n$ is the
closed convex hull of the same printed exact-rank set.
\footnote{Source: \path{Notes/WolfNoteTexSource/ch03_positive_not_completely.tex},
lines 229--243. Tracked by
\href{https://github.com/LionSR/QICLean/issues/23}{QICLean issue \#23}.}

This exact-rank wording again fails without the bound $n\le\min\{d,d'\}$.
If $n>\min\{d,d'\}$, no vector has Schmidt rank exactly $n$, so the witness
condition is vacuous. Taking $W=-\one$ then makes $\tr(W\rho)=-1<0$ for every
density operator, whereas no such $\rho$ can have Schmidt number at least
$n+1$. Thus the printed equivalence is false in precisely the range where its
test set is empty.

The dimension-generic correction is
\begin{align}
    \operatorname{sn}(\rho)>n
    \quad\Longleftrightarrow\quad
    \exists W=W^\dagger:\quad
    \tr(W\rho)<0\quad\text{and}\quad
    \langle\psi,W\psi\rangle&\ge0
    &&\text{whenever }\operatorname{SR}(\psi)\le n.
    \tag{corrected Wolf 3.3}
\end{align}
This is exactly the separating-hyperplane statement for $S_n$: each pure
projector in a defining decomposition of $S_n$ has Schmidt rank at most $n$.
It is formalized by
\leanid{Matrix.not_hasSchmidtNumberLE_iff_exists_witness} in
\doclink{QICLean/Channel/EntanglementWitness} and is the witness input to the
rectangular positive-map characterization
\leanid{Matrix.hasSchmidtNumberLE_iff_forall_isNPositiveMap_tensorMapId_posSemidef}
in \doclink{QICLean/Channel/PositiveMapDetection}.

When $n\le\min\{d,d'\}$, the exact-rank witness test can again be recovered
from the bounded-rank test by the same density-and-continuity argument. That
additional argument is not the separating-hyperplane proof printed in
Proposition~3.3 and is not used by the formal theorem.

\section{Conclusion}

Wolf's exact-rank wording in Proposition~3.1 is false under the printed
hypothesis $k\le d$, and the same wording in Proposition~3.3 is false when
$n>\min\{d,d'\}$. Replacing exact Schmidt rank by Schmidt rank at most the
given bound yields the dimension-generic Choi and witness criteria and matches
the proofs through compressions and the convex set $S_n$. The formal statements
and the Blueprint use this corrected wording; no additional route is attributed
to the source.

\bibliographystyle{alpha}
\bibliography{references}

\end{document}
